\section{Related Works}

\subsection{LLM for Planning}\label{sec:related-planning}
Large language models (LLMs), enriched with diverse world knowledge and human priors during their pretraining phase,  possess the ability to understanding and reasoning\cite{yao2023react, guo2025deepseek, yang2024qwen2, yan2023codescope, yi2024survey}.
When provided with relevant task information, LLMs are able to decompose complex, long-term tasks into simpler, short-term tasks, each focused on specific objectives\cite{khot2023decomposedpromptingmodularapproach, yuan2023tasklamaprobingcomplextask}.
Numerous approaches have leveraged Task and Motion Planning (TAMP) to realize such task decomposition in robotics scenes\cite{chen2023nl2tl, chen2024autotamp}.
However, when using RL for low-level control strategies, a significant gap arises between high-level planning and low-level control—how to efficiently leverage the natural language plans generated by LLMs to guide policy training?
SayCan\cite{brohan2023can} addresses this challenge by utilizing a set of pretrained policies, where the LLM simply selects the appropriate policy based on the plans. Plan-Seq-Learn\cite{dalal2024plan} uses traditional motion planning to link subtasks, alleviating the burden on RL policies.
In text-based environments, ELLM\cite{du2023guiding} trains RL policies by describing state transition in natural language and computing the similarity between the executed actions and the plan, using this similarity as a reward signal. However, this approach is computationally expensive and not suitable for robotics tasks.
In contrast, our method constructs an efficient and dense reward signal through the Progress Reward Model (PRM), effectively bridging the gap between high-level planning and low-level policy training.

\subsection{LLM for Reward Generation}\label{sec:related-rewards}
Reward functions are the primary mechanism to guide RL training\cite{devidze2024informativenessrewardfunctionsreinforcement, eschmann2021reward}, leading to numerous approaches that utilize LLMs to generate executable reward functions that provide dense rewards. Text2Reward\cite{xie2024text2reward} uses LLMs to generate executable python functions based on a goal described in natural language, with human involvement in the loop.
Eureka\cite{ma2023eureka} extends this paradigm by introducing evolutionary search and automatic reward refinement through LLMs.
REvolve\cite{hazra2024revolve} translates implicit human knowledge into explicit reward functions for training RL with human feedback.
However, these methods treat the task as a whole and do not account for the potential stage-wise features of complex tasks.
As a result, the generation process of reward function tend to be highly complex and often require human feedback\cite{xie2024text2reward, hazra2024revolve} or extensive search\cite{ma2023eureka}.
In contrast, our approach simplifies the reward generation process by decomposing the task into short-term, simple subtasks, providing more specific and clear guidance for each stage of the task.
